% Section fragment; the enclosing report supplies theorem environments.
\section{The 215-state composition}
\label{sec:composition215}

The primary new table has 215 working states and 430 transition rows. It combines
changes from two credited public 216-state predecessors: the partial-selector
source layout and native 164/183 fold; the MP-only read schedule and 210 command
permutation; and the contextual replacement of native row 81/0. Its additional
contribution is an exact allocation in which these changes coexist. The public
predecessors are preserved as unchanged, hash-pinned archives, with their
original reports and attribution.

\subsection{Source construction and logical correspondence}

Start with the public partial-selector source, whose complete guarded ordinal
graph is equal to that of K217. In its original primitive ordinals, delete
$[23,37)$ and $[187,201)$, and replace each of $[228,254)$ and $[4659,4685)$ by the
local counter read followed by the existing counter load. The 14-leaf read
factors as $(W_{11},T_3)$, and the load has 26 leaves. Use
\[
  (W_{11},(T_3,L_{26}))
\]
for both replacements. Each of the three factors is closed and halt-free, so
this association preserves every effect and guarded successor. It also keeps
the maximum consecutive zero run at 6.

Permute the three closed command bodies at paths
\texttt{111}, \texttt{110101}, and \texttt{110100101} in order 210. Tags 0, 1, 2 select
NIP, DROP, MP; a selector at least 3 still selects AX with residual selector $n-3$.
The resulting source has 202 distinct subtrees, 4751 primitive occurrences,
maximum depth 47, and maximum zero run 6.

\begin{lemma}[Complete source correspondence]
The composed 210 source and the public calibrated 210 source induce the same
guarded transition system on primitive ordinals and all nine natural-number
registers.
\end{lemma}
\begin{proof}
The finite certificate compares their entire ordinal-indexed graphs. For an
increment it records the register and successor; for a decrement it records the
register and both the zero and positive successors; for break or continue it
records the resolved identity successor; and for halt it records termination.
An ancestor-interval algorithm and an independent binary-address arithmetic
algorithm agree, including the no-stray-suffix condition. An additional source
auditor checks the command selector partition, exact wire fields, and each
normal return route. These graphs are literally equal for the two sources.
Induction on abstract primitive steps therefore gives the identity
correspondence on registers and ordinals, including termination, for every
initial register vector. No formula-decoding assumption or experimental
numeric bound is used.
\end{proof}

The numerical-calculus and command-calibration theorems of the public 210
construction therefore apply unchanged. The new blank initialization reaches
the same logical vector $(4,0,0,0,0,0,0,0,0)$ as that construction.

\subsection{Constrained physical allocation}

Of the new source's 179 internal subtree nodes, 157 retain old roles represented
by 143 physical decoder states. The other 22 nodes are assigned injectively to a
23-state pool consisting of retired decoder states and the dispatcher root.
The assignment leaves only native 185 unused. The 50 remaining engine states,
143 retained decoder states, and 22 assigned decoder states give
\[
  50+143+22=215.
\]
This is a count of the actual complete table: each working state has both
binary transition rows, and 254 is its sole native external halt.

Rows $(11,0),(26,0),(26,1),(37,0),(92,0),(213,0),(213,1)$ in the pool continue to
serve arithmetic or marker helpers and are preserved exactly. The public local
replacement $81/0=0L111$ frees row 92/1. Its original proof compares the two
possible symbols immediately to the left of the head and gives an exact
three-step whole-tape equality. Both inherited marker folds retain their
literal provider fragments: the 69 fold uses 213, 92 and 90, while the 164 fold uses
201, 195 and 196. Fresh independent local checks also verify the folds after
unfolding them to their immediate predecessors.

Relative to the public partial-selector 216 table, the new table changes 36
retained rows, deletes both rows of 185, and introduces no state. The complete
before/after recipe is \texttt{recipes/table210.json}. Its fixed assignment
identifies every new node by source path, so reconstruction requires no
optimization or search. The standard-library builder checks all internal
prefixes against that assignment. The primary native table hash is
\begin{center}\small\ttfamily
04b4dbb537306f40f8965ad6a0ee6034e9953f079c233153e9696acdf06c2712.
\end{center}

\subsection{Endpoint proof}

The finite source and allocation certificates are supplemented by the full
whole-tape symbolic runtime. All 6176 guarded primitive cases are proved over
all natural primary-register values, arbitrary natural extra gap, a nonempty
finite spare list, and zero exterior. Every primitive block has positive finite
clock. The independent inventory covers all 430 rows, with 1698790 ordered rule
applications and minimum primitive clock 6. This supplies both directions of
the source-to-machine simulation and excludes an additional halting exit.

The complete blank prefix has 95726 transitions. Independent sparse and dense
replays agree on every transition and the whole final tape. The endpoint is
native dispatcher 84, head $-8$, extra gap 1, the register vector
$(4,0,0,0,0,0,0,0,0)$, and 40 spare blocks of value 1. There are 95 one cells and the
tape is zero everywhere else. Combining this initialization with the quantified
runtime, the source-correspondence lemma, and the inherited calibration gives
the ZFC-provable equivalence between nonhalting of the exact 215-state table and
$\operatorname{Con}(\mathrm{ZFC})$.

These are ordinary mathematical certificates formalizable in ZFC. No new
proof-assistant-kernel endpoint, absolute consistency proof, or minimality
among all equivalent machines is asserted.
